\documentclass[12pt]{article}
\usepackage[utf8]{inputenc}
\usepackage{listings}
\usepackage{xcolor}
\usepackage{geometry}
\usepackage{parskip}

\geometry{margin=2.5cm}

\definecolor{codebg}{RGB}{245,245,245}
\definecolor{keyword}{RGB}{0,112,0}
\definecolor{comment}{RGB}{128,128,128}
\definecolor{string}{RGB}{180,0,0}

\lstdefinestyle{pythonstyle}{
    backgroundcolor=\color{codebg},
    basicstyle=\ttfamily\small,
    keywordstyle=\color{keyword}\bfseries,
    stringstyle=\color{string},
    commentstyle=\color{comment}\itshape,
    language=Python,
    frame=single,
    rulecolor=\color{gray!40},
    breaklines=true,
    showstringspaces=false,
    tabsize=4,
    captionpos=b,
}

\title{\textbf{Wind Turbine Blade Fault Classification\\Using a Simple CNN}\\[0.5em]
\large Step 7 --- Build the CNN Model}
\author{Amreen Batool}
\date{}

\begin{document}
\maketitle

\section*{Step 7 --- Build the CNN Model}

\begin{lstlisting}[style=pythonstyle]
model = models.Sequential([
    layers.Input(shape=(128, 128, 3)),

    layers.Conv2D(32, (3, 3), activation='relu'),
    layers.MaxPooling2D(2, 2),

    layers.Conv2D(64, (3, 3), activation='relu'),
    layers.MaxPooling2D(2, 2),

    layers.Conv2D(128, (3, 3), activation='relu'),
    layers.MaxPooling2D(2, 2),

    layers.Flatten(),
    layers.Dense(64, activation='relu'),
    layers.Dropout(0.5),
    layers.Dense(1, activation='sigmoid')
])
\end{lstlisting}

\subsection*{Explanation of Each Layer}

\begin{itemize}
    \item \texttt{layers.Input(shape=(128, 128, 3))} \\
          Input image size is $128 \times 128$ pixels with 3 colour channels (RGB).

    \item \texttt{Conv2D(32, (3,3), activation='relu')} \\
          First convolution layer. Detects basic image patterns like edges and textures using 32 filters.

    \item \texttt{MaxPooling2D(2,2)} \\
          Reduces the feature map size by half, keeping only the most important information.

    \item \textbf{Second convolution block} (\texttt{Conv2D(64, ...)}) \\
          Learns more detailed features using 64 filters.

    \item \textbf{Third convolution block} (\texttt{Conv2D(128, ...)}) \\
          Learns deeper and more complex features using 128 filters.

    \item \texttt{Flatten()} \\
          Converts the 2D feature maps into a 1D vector so they can be passed to the dense layers.

    \item \texttt{Dense(64, activation='relu')} \\
          Fully connected layer for learning final high-level patterns.

    \item \texttt{Dropout(0.5)} \\
          Randomly sets 50\% of neurons to zero during training. Helps reduce overfitting.

    \item \texttt{Dense(1, activation='sigmoid')} \\
          Final output layer for binary classification. Produces a probability between 0 and 1. \\
          Values $> 0.5$ are classified as \texttt{Healthy}; values $\leq 0.5$ as \texttt{Faulty}.
\end{itemize}

\end{document}
